\documentclass[11pt]{article}

\usepackage[utf8]{inputenc}
\usepackage[T1]{fontenc}
\usepackage[spanish,es-nodecimaldot]{babel}
\usepackage{amsmath,amssymb,amsthm}
\usepackage{mathtools}
\usepackage{graphicx}
\usepackage{booktabs}
\usepackage{geometry}
\usepackage{hyperref}
\usepackage{xcolor}
\usepackage{enumitem}
\usepackage{siunitx}
\usepackage{array}

\geometry{margin=2.5cm}
\hypersetup{
    colorlinks=true,
    linkcolor=blue!50!black,
    urlcolor=blue!50!black
}

\title{Adaptaci\'on del Proyecto a lo Visto en Clase\\Paso a Paso y con Justificaci\'on Matem\'atica}
\author{Plataforma Antivibratoria Activa}
\date{14 de abril de 2026}

\begin{document}

\maketitle
\tableofcontents
\newpage

\section{Prop\'osito de este documento}

La idea de este texto es tomar el proyecto que ya planteamos y \textbf{reescribirlo exactamente en la forma en que se conecta con lo que has visto en clase}. No est\'a pensado como un informe corto, sino como una explicaci\'on muy detallada para que entiendas:

\begin{itemize}[leftmargin=1.2cm]
\item qu\'e hicimos en cada paso;
\item por qu\'e ese paso tiene sentido matem\'aticamente;
\item c\'omo se conecta con modelado no lineal, linealizaci\'on, estabilidad, respuesta escal\'on, PID e IMC;
\item c\'omo defenderlo si te preguntan en clase o en la sustentaci\'on.
\end{itemize}

El hilo matem\'atico que seguimos en este documento es el mismo que aparece en las diapositivas del curso:

\begin{enumerate}[leftmargin=1.2cm]
\item obtener el modelo din\'amico;
\item expresarlo en variables de estado;
\item linealizar alrededor de un punto de equilibrio;
\item obtener una funci\'on de transferencia para el caso SISO;
\item estudiar estabilidad del modelo aproximado;
\item analizar la respuesta temporal;
\item dise\~nar un controlador usando lo visto hasta IMC;
\item comparar el modelo realista y el aproximado.
\end{enumerate}

\section{Qu\'e sistema escogimos y por qu\'e}

Escogimos una \textbf{plataforma antivibratoria activa}. F\'isicamente, esto representa una mesa o plataforma donde se monta un equipo sensible. La base de esa plataforma vibra por una perturbaci\'on externa, y lo que queremos es que la parte superior donde est\'a el equipo se mueva lo menos posible.

La aplicaci\'on es totalmente defendible porque el aislamiento vibratorio aparece en:

\begin{itemize}[leftmargin=1.2cm]
\item mesas \'opticas;
\item soportes de sensores;
\item plataformas de precisi\'on;
\item protecci\'on de instrumentaci\'on delicada;
\item sistemas de rechazo de vibraci\'on industrial.
\end{itemize}

\subsection{Por qu\'e esta idea cumple el enunciado}

La actividad pide un sistema:

\begin{itemize}[leftmargin=1.2cm]
\item SISO;
\item de al menos tercer orden;
\item con un modelo realista;
\item con un modelo aproximado o linealizado;
\item con an\'alisis de estabilidad;
\item con dise\~no de controlador;
\item con evaluaci\'on de seguimiento y rechazo a perturbaciones.
\end{itemize}

Nuestro proyecto cumple cada punto:

\begin{itemize}[leftmargin=1.2cm]
\item \textbf{SISO}: la entrada es la se\~nal al actuador $u(t)$ y la salida es la posici\'on de la plataforma $x(t)$.
\item \textbf{Tercer orden}: la planta tiene dos estados mec\'anicos y un estado adicional del actuador.
\item \textbf{Modelo realista}: incluye no linealidades y restricciones f\'isicas.
\item \textbf{Modelo aproximado}: se obtiene linealizando alrededor del equilibrio.
\item \textbf{Control}: se dise\~na mediante IMC.
\item \textbf{Evaluaci\'on}: se comparan ambos modelos en lazo abierto y lazo cerrado.
\end{itemize}

\section{Paso 1: Modelo f\'isico no lineal}

Lo primero que se hace, exactamente como se explica en las diapositivas de estabilidad, es partir de una representaci\'on general no lineal:

\begin{equation}
\dot{x}=f(x,u), \qquad y=h(x,u).
\end{equation}

Eso significa que antes de hablar de funci\'on de transferencia o de PID, lo correcto es construir el sistema como \textbf{modelo din\'amico no lineal}.

\subsection{Variables f\'isicas}

Definimos:

\begin{itemize}[leftmargin=1.2cm]
\item $x(t)$: desplazamiento absoluto de la plataforma.
\item $\dot{x}(t)$: velocidad de la plataforma.
\item $\ddot{x}(t)$: aceleraci\'on de la plataforma.
\item $z(t)$: desplazamiento de la base vibrante.
\item $\dot{z}(t)$: velocidad de la base vibrante.
\item $u(t)$: se\~nal de control aplicada al actuador.
\item $F_a(t)$: fuerza generada por el actuador.
\end{itemize}

\subsection{Par\'ametros f\'isicos}

Los par\'ametros usados fueron:

\begin{center}
\begin{tabular}{ccc}
\toprule
Par\'ametro & Valor & Significado \\
\midrule
$m$ & $12$ & masa de la plataforma [kg] \\
$k$ & $1800$ & rigidez lineal [N/m] \\
$c$ & $95$ & amortiguamiento viscoso [N\,s/m] \\
$k_3$ & $6\times 10^4$ & rigidez c\'ubica [N/m$^3$] \\
$F_c$ & $1$ & fricci\'on seca suavizada [N] \\
$K_a$ & $220$ & ganancia del actuador [N/V] \\
$\tau_a$ & $0.04$ & constante de tiempo del actuador [s] \\
\bottomrule
\end{tabular}
\end{center}

\subsection{Ecuaci\'on mec\'anica}

Aplicamos la segunda ley de Newton:
\[
\sum F = m\ddot{x}.
\]

La masa recibe:

\begin{itemize}[leftmargin=1.2cm]
\item una fuerza de actuaci\'on $F_a$;
\item una fuerza de suspensi\'on que depende del movimiento relativo con la base;
\item una fuerza de tope mec\'anico si la masa sobrepasa el recorrido admisible.
\end{itemize}

Por tanto:
\begin{equation}
m\ddot{x} = F_a - F_{\text{susp}} - F_{\text{stop}}.
\label{eq:newton}
\end{equation}

\subsection{Por qu\'e la suspensi\'on depende de $x-z$ y $\dot{x}-\dot{z}$}

Esto es muy importante conceptualmente. Si la base se mueve, el resorte y el amortiguador no dependen de la posici\'on absoluta de la plataforma sino de la \textbf{deformaci\'on relativa}. Entonces:

\begin{itemize}[leftmargin=1.2cm]
\item deformaci\'on del resorte: $x-z$;
\item velocidad relativa: $\dot{x}-\dot{z}$.
\end{itemize}

Si la base y la plataforma se mueven exactamente igual, la suspensi\'on no se est\'a deformando ni disipando energ\'ia relativa. Esa es la raz\'on f\'isica de usar variables relativas.

\subsection{Fuerza de suspensi\'on realista}

La definimos como:
\begin{equation}
F_{\text{susp}} =
k(x-z)
c(\dot{x}-\dot{z})
k_3(x-z)^3
F_c\tanh\left(\frac{\dot{x}-\dot{z}}{v_\varepsilon}\right).
\label{eq:fsusp}
\end{equation}

Cada t\'ermino tiene una justificaci\'on:

\begin{itemize}[leftmargin=1.2cm]
\item $k(x-z)$: elasticidad lineal.
\item $c(\dot{x}-\dot{z})$: amortiguamiento lineal.
\item $k_3(x-z)^3$: no linealidad de rigidez.
\item $F_c\tanh((\dot{x}-\dot{z})/v_\varepsilon)$: fricci\'on seca suavizada.
\end{itemize}

\subsection{Por qu\'e agregamos no linealidades}

Porque el enunciado exige que el modelo realista tenga efectos f\'isicos que no est\'en en el modelo aproximado. Si no se hace eso, ambos modelos ser\'ian casi el mismo sistema y se perder\'ia el objetivo central de la actividad.

\subsection{Din\'amica del actuador}

No modelamos el actuador como fuerza instant\'anea. Hicimos:
\begin{equation}
\tau_a\dot{F}_a + F_a = K_a\,\mathrm{sat}(u).
\label{eq:actuador}
\end{equation}

Esto quiere decir:

\begin{itemize}[leftmargin=1.2cm]
\item si aplicas un voltaje o mando $u$, la fuerza no aparece de inmediato;
\item el actuador tiene una din\'amica de primer orden;
\item adem\'as, la se\~nal tiene saturaci\'on porque un actuador real no puede empujar infinitamente.
\end{itemize}

\subsection{Topes mec\'anicos}

Agregamos topes mec\'anicos para que el desplazamiento de la plataforma no sea f\'isicamente ilimitado. Esto hace m\'as realista el modelo y ayuda a justificar que la planta tiene restricciones reales.

\section{Paso 2: Llevar el modelo no lineal a espacio de estado}

Siguiendo exactamente la secuencia matem\'atica de las diapositivas, el siguiente paso es expresar el sistema en variables de estado.

\subsection{Elecci\'on de estados}

Escogemos:
\begin{equation}
x_1 = x, \qquad x_2 = \dot{x}, \qquad x_3 = F_a.
\end{equation}

Esta elecci\'on es natural porque:

\begin{itemize}[leftmargin=1.2cm]
\item $x_1$ representa la posici\'on;
\item $x_2$ representa la velocidad;
\item $x_3$ representa el estado interno del actuador.
\end{itemize}

\subsection{Ecuaciones de estado}

Entonces:

\begin{equation}
\dot{x}_1 = x_2,
\end{equation}

\begin{equation}
\dot{x}_2 =
\frac{1}{m}\Big(
x_3 - F_{\text{susp}} - F_{\text{stop}}
\Big),
\end{equation}

\begin{equation}
\dot{x}_3 =
\frac{-x_3 + K_a\,\mathrm{sat}(u)}{\tau_a}.
\end{equation}

Y la salida es:
\begin{equation}
y = x_1.
\end{equation}

Con esto ya tenemos formalmente una representaci\'on del tipo
\[
\dot{x}=f(x,u), \qquad y=h(x,u),
\]
que es precisamente lo que se muestra en tus diapositivas de estabilidad como punto de partida.

\section{Paso 3: Encontrar el punto de equilibrio}

Antes de linealizar, se necesita un punto de equilibrio.

\subsection{Qu\'e significa equilibrio}

Un punto de equilibrio es un punto donde, si el sistema empieza all\'i y no hay perturbaciones adicionales, entonces permanece all\'i. En otras palabras:
\[
\dot{x}=0.
\]

\subsection{Equilibrio escogido}

Escogimos el equilibrio natural:
\[
x^\star = 0,\qquad
\dot{x}^\star = 0,\qquad
z^\star = 0,\qquad
\dot{z}^\star = 0,\qquad
F_a^\star = 0,\qquad
u^\star = 0.
\]

Este equilibrio tiene sentido porque representa la plataforma en reposo, sin excitaci\'on y sin desplazamiento.

\section{Paso 4: Linealizaci\'on}

Ahora seguimos exactamente el esquema de tus clases:

\begin{equation}
\dot{\hat{x}} = A\hat{x} + B\hat{u}, \qquad \hat{y}=C\hat{x}+D\hat{u}.
\end{equation}

La idea es aproximar el sistema no lineal alrededor del equilibrio usando Jacobianos.

\subsection{Qu\'e se desprecia en la linealizaci\'on}

Alrededor del equilibrio, los t\'erminos no lineales se pueden aproximar o despreciar:

\begin{itemize}[leftmargin=1.2cm]
\item la rigidez c\'ubica desaparece en primera aproximaci\'on porque es de orden superior en la deformaci\'on;
\item la fricci\'on seca suavizada se aproxima cerca del equilibrio, pero para el modelo de dise\~no la ignoramos;
\item los topes mec\'anicos no act\'uan mientras estemos cerca del origen;
\item la saturaci\'on se ignora en el modelo nominal para dise\~no.
\end{itemize}

\subsection{Modelo lineal resultante}

Con la base fija en cero para estudiar la planta respecto a la entrada de control:
\begin{equation}
m\ddot{x} + c\dot{x} + kx = F_a,
\label{eq:lineal_mecanico}
\end{equation}
\begin{equation}
\tau_a\dot{F}_a + F_a = K_a u.
\label{eq:lineal_act}
\end{equation}

\subsection{Espacio de estado lineal}

Con los estados escogidos:
\[
\hat{x}=
\begin{bmatrix}
\hat{x}_1\\
\hat{x}_2\\
\hat{x}_3
\end{bmatrix}

=
\begin{bmatrix}
\hat{x}\\
\widehat{\dot{x}}\\
\widehat{F_a}
\end{bmatrix},
\]
obtenemos:
\begin{equation}
\dot{\hat{x}}=
\begin{bmatrix}
0 & 1 & 0\\
-\frac{k}{m} & -\frac{c}{m} & \frac{1}{m}\\
0 & 0 & -\frac{1}{\tau_a}
\end{bmatrix}
\hat{x}

+
\begin{bmatrix}
0\\
0\\
\frac{K_a}{\tau_a}
\end{bmatrix}
\hat{u},
\end{equation}

\begin{equation}
\hat{y}=
\begin{bmatrix}
1 & 0 & 0
\end{bmatrix}
\hat{x}.
\end{equation}

Entonces:
\[
A=
\begin{bmatrix}
0 & 1 & 0\\
-\frac{k}{m} & -\frac{c}{m} & \frac{1}{m}\\
0 & 0 & -\frac{1}{\tau_a}
\end{bmatrix},
\quad
B=
\begin{bmatrix}
0\\
0\\
\frac{K_a}{\tau_a}
\end{bmatrix},
\quad
C=
\begin{bmatrix}
1 & 0 & 0
\end{bmatrix},
\quad
D=0.
\]

\section{Paso 5: Obtener la funci\'on de transferencia}

En tus diapositivas aparece que, para el caso SISO, la funci\'on de transferencia puede obtenerse como:
\[
G(s)=C(sI-A)^{-1}B + D.
\]

En este proyecto, sin embargo, es incluso m\'as f\'acil obtenerla directamente desde las ecuaciones diferenciales.

Aplicando Laplace a \eqref{eq:lineal_mecanico} y \eqref{eq:lineal_act}:
\[
(ms^2+cs+k)X(s)=F_a(s),
\]
\[
(\tau_as+1)F_a(s)=K_aU(s).
\]

Eliminando $F_a(s)$:
\begin{equation}
G(s)=\frac{X(s)}{U(s)}
=
\frac{K_a}{(\tau_as+1)(ms^2+cs+k)}.
\label{eq:G}
\end{equation}

\subsection{Por qu\'e esta funci\'on de transferencia es clave}

Porque esta es la planta \textbf{aproximada} sobre la que dise\~namos el controlador. Esta planta es suficientemente simple para an\'alisis y suficientemente fiel para capturar lo esencial de la din\'amica.

\section{Paso 6: Verificar que el sistema es de tercer orden}

Este paso es simple pero importante, porque el enunciado lo exige.

El denominador de \eqref{eq:G} es:
\[
(\tau_as+1)(ms^2+cs+k).
\]

Ese denominador tiene grado tres. Por tanto:
\[
\deg D(s)=3.
\]

Eso significa que la planta aproximada es de tercer orden.

\section{Paso 7: An\'alisis de estabilidad}

Seg\'un lo visto en clase, una forma central de estudiar estabilidad en sistemas LTI es observar los polos del sistema o, equivalentemente, aplicar criterios como Routh-Hurwitz.

\subsection{Denominador de la planta}

Con los valores num\'ericos:
\[
\tau_a = 0.04,\quad
m = 12,\quad
c=95,\quad
k=1800,
\]
el denominador queda:
\begin{equation}
D(s)=0.48s^3 + 15.8s^2 + 167s + 1800.
\label{eq:den}
\end{equation}

\subsection{Interpretaci\'on con polos}

Los polos calculados son aproximadamente:
\[
s_1=-25,\qquad
s_{2,3}=-3.958\pm j\,11.590.
\]

Todos est\'an en el semiplano izquierdo, luego la planta aproximada es estable.

\subsection{Interpretaci\'on con Routh-Hurwitz}

Como en clase les explicaron RH, tambi\'en es muy natural defender estabilidad as\'i.

Para un polinomio c\'ubico:
\[
a_3s^3+a_2s^2+a_1s+a_0,
\]
una condici\'on necesaria y suficiente de estabilidad es:

\begin{itemize}[leftmargin=1.2cm]
\item $a_3>0$,
\item $a_2>0$,
\item $a_1>0$,
\item $a_0>0$,
\item $a_2a_1>a_3a_0$.
\end{itemize}

En nuestro caso:
\[
a_3=0.48,\quad
a_2=15.8,\quad
a_1=167,\quad
a_0=1800.
\]

Todos son positivos. Adem\'as:
\[
a_2a_1 = 15.8\cdot167 = 2638.6,
\]
\[
a_3a_0 = 0.48\cdot1800 = 864.
\]

Como
\[
2638.6 > 864,
\]
entonces por Routh-Hurwitz la planta aproximada es estable.

\subsection{Por qu\'e este paso importa}

Porque no basta con decir ``el modelo se ve razonable''. El curso les exige justificar estabilidad con herramientas formales, y RH es exactamente una de las herramientas que viste en las diapositivas.

\section{Paso 8: An\'alisis temporal de la planta}

De acuerdo con lo que aparece en las diapositivas de respuesta escal\'on y sintonizaci\'on, una vez obtenida la planta se deben estudiar m\'etricas como:

\begin{itemize}[leftmargin=1.2cm]
\item tiempo de subida;
\item tiempo pico;
\item tiempo de establecimiento;
\item sobrepico;
\item error en estado estacionario.
\end{itemize}

\subsection{Por qu\'e esto se hace antes del controlador}

Porque primero debemos entender c\'omo se comporta la planta sola. Eso nos dice si es lenta, r\'apida, oscilatoria, con alto sobrepico, etc.

\subsection{Interpretaci\'on de nuestra planta}

La planta tiene:

\begin{itemize}[leftmargin=1.2cm]
\item un polo real r\'apido del actuador;
\item un par dominante complejo conjugado asociado a la parte masa-resorte-amortiguador.
\end{itemize}

Eso significa que el comportamiento temporal va a estar dominado por una din\'amica oscilatoria amortiguada, con un modo r\'apido adicional que se apaga pronto.

\section{Paso 9: Comparaci\'on entre modelo realista y aproximado}

Este paso es central en la actividad.

\subsection{Qu\'e hicimos exactamente}

Tomamos:

\begin{itemize}[leftmargin=1.2cm]
\item el modelo aproximado lineal;
\item el modelo realista no lineal;
\item la misma entrada de prueba;
\item y comparamos sus respuestas.
\end{itemize}

En lazo abierto se us\'o un escal\'on peque\~no en el actuador. La raz\'on para usar una entrada peque\~na es que, cerca del equilibrio, el modelo lineal deber\'ia aproximar bien al sistema realista.

\subsection{Qu\'e se espera observar}

\begin{itemize}[leftmargin=1.2cm]
\item Si la entrada es peque\~na, ambos modelos deben parecerse razonablemente.
\item El modelo realista puede mostrar menos sobrepico o cambios en el valor final por la fricci\'on y las no linealidades.
\item Cuando la excitaci\'on crece, la diferencia entre ambos modelos se vuelve m\'as notable.
\end{itemize}

\section{Paso 10: Elecci\'on del controlador}

Como en tus archivos aparece contenido tanto de PID como de IMC, lo correcto es explicar por qu\'e elegimos \textbf{IMC} y no una sinton\'ia PID cl\'asica.

\subsection{Raz\'on 1: el enunciado lo permite}

La actividad dice expl\'icitamente que se puede dise\~nar el controlador con los temas vistos hasta IMC.

\subsection{Raz\'on 2: la planta aproximada es buena candidata}

La planta aproximada:

\begin{itemize}[leftmargin=1.2cm]
\item es estable;
\item no tiene tiempos muertos expl\'icitos;
\item no tiene ceros no m\'inimos.
\end{itemize}

Eso hace a IMC una metodolog\'ia muy natural para este caso.

\subsection{Raz\'on 3: IMC deja una matem\'atica muy limpia}

Con IMC se puede imponer directamente una din\'amica nominal deseada en lazo cerrado. Esto facilita mucho la explicaci\'on y la justificaci\'on del dise\~no.

\section{Paso 11: Dise\~no IMC paso a paso}

\subsection{Planta nominal}

Tomamos:
\[
\tilde{G}_p(s)=\frac{K_a}{(\tau_as+1)(ms^2+cs+k)}.
\]

\subsection{Filtro IMC}

Escogemos:
\begin{equation}
F(s)=\frac{1}{(\lambda s+1)^3}.
\label{eq:filtro}
\end{equation}

\subsection{Por qu\'e se escoge un filtro de orden 3}

Porque la inversa de la planta nominal tiene grado relativo negativo. Si no se a\~nade un filtro de orden suficiente, el controlador resultante no ser\'ia propio ni realizable.

Como la planta es de tercer orden, una elecci\'on natural es un filtro de tercer orden.

\subsection{Controlador interno}

\begin{equation}
Q(s)=\tilde{G}_p^{-1}(s)F(s)
=
\frac{(\tau_as+1)(ms^2+cs+k)}{K_a(\lambda s+1)^3}.
\end{equation}

\subsection{Controlador equivalente}

Usamos la expresi\'on vista en IMC:
\begin{equation}
G_c(s)=\frac{Q(s)}{1-\tilde{G}_p(s)Q(s)}.
\end{equation}

Como
\[
\tilde{G}_p(s)Q(s)=\frac{1}{(\lambda s+1)^3},
\]
resulta:
\begin{equation}
G_c(s)=
\frac{(\tau_as+1)(ms^2+cs+k)}
{K_a\lambda s(\lambda^2s^2+3\lambda s+3)}.
\label{eq:gc}
\end{equation}

\subsection{Qu\'e significa el integrador}

En \eqref{eq:gc} aparece un factor $1/s$. Eso es muy importante porque significa que el controlador tiene acci\'on integral, lo cual ayuda a reducir error en estado estacionario.

\subsection{Qu\'e papel juega $\lambda$}

$\lambda$ regula la rapidez deseada del sistema nominal:

\begin{itemize}[leftmargin=1.2cm]
\item menor $\lambda$: respuesta m\'as r\'apida, m\'as agresiva;
\item mayor $\lambda$: respuesta m\'as lenta, m\'as robusta.
\end{itemize}

En el proyecto usamos:
\[
\lambda = 0.10\ \text{s}.
\]

\subsection{Din\'amica nominal esperada}

Con IMC, la transferencia nominal en lazo cerrado queda:
\[
T(s)=\frac{1}{(\lambda s+1)^3}.
\]

Eso implica una respuesta suave, sin sobrepico nominal ideal, y con un comportamiento controlado por $\lambda$.

\section{Paso 12: Relaci\'on con lo visto de PID}

Aunque el controlador dise\~nado no se forz\'o a ser un PID ideal cl\'asico, s\'i est\'a conectado con lo visto en clase por varias razones:

\begin{itemize}[leftmargin=1.2cm]
\item contiene acci\'on integral;
\item tiene una estructura racional que puede entenderse como compensador din\'amico de orden mayor;
\item IMC y PID aparecen juntos en tus materiales, precisamente porque muchas veces IMC lleva a controladores tipo PID o parecidos a PID cuando la planta tiene cierta estructura.
\end{itemize}

Si te preguntan por qu\'e no se escogi\'o un PID puro, una respuesta honesta y matem\'aticamente s\'olida es:

\begin{quote}
Se prefiri\'o mantener el dise\~no directo por IMC porque la planta es de tercer orden y no se quiso forzar una reducci\'on adicional de la din\'amica solo para obtener un PID ideal. Como el curso permite IMC y el objetivo es comparar modelos, esta elecci\'on es consistente con el enunciado.
\end{quote}

\section{Paso 13: Simulaci\'on en lazo cerrado}

Despu\'es de dise\~nar el controlador, hicimos lo que pide el enunciado:

\begin{itemize}[leftmargin=1.2cm]
\item aplicar el controlador al modelo aproximado;
\item aplicar el mismo controlador al modelo realista;
\item estudiar seguimiento a referencia variable;
\item estudiar rechazo a perturbaciones.
\end{itemize}

\subsection{Referencia variable}

No dejamos una sola referencia constante. Se usaron varios escalones:

\begin{itemize}[leftmargin=1.2cm]
\item $0$ m al inicio;
\item $0.004$ m;
\item $-0.0025$ m;
\item $0.003$ m.
\end{itemize}

Esto permite mostrar que el controlador no solo funciona para un \'unico cambio, sino para varios cambios sucesivos.

\subsection{Perturbaci\'on de base}

La perturbaci\'on se aplic\'o directamente a la base. Esto tambi\'en se alinea bien con la interpretaci\'on f\'isica del sistema:

\begin{itemize}[leftmargin=1.2cm]
\item primero no hay perturbaci\'on;
\item luego aparece una vibraci\'on senoidal dominante;
\item despu\'es una mezcla de senoidales m\'as peque\~nas.
\end{itemize}

\section{Paso 14: M\'etricas de desempe\~no}

Tal como lo pide el enunciado, las m\'etricas a evaluar son:

\begin{itemize}[leftmargin=1.2cm]
\item tiempo de establecimiento;
\item tiempo de subida;
\item overshoot;
\item error en estado estacionario.
\end{itemize}

\subsection{Por qu\'e estas m\'etricas importan}

\begin{itemize}[leftmargin=1.2cm]
\item El \textbf{tiempo de subida} dice qu\'e tan r\'apido reacciona el sistema.
\item El \textbf{tiempo de establecimiento} dice cu\'anto tarda en ``quedar quieto'' cerca del valor final.
\item El \textbf{overshoot} dice cu\'anto se pasa.
\item El \textbf{error estacionario} dice si realmente llega al valor deseado.
\end{itemize}

\subsection{M\'etrica adicional}

Adem\'as, introdujimos el error RMS en la ventana de perturbaci\'on, porque para un sistema antivibratorio esa m\'etrica es muy informativa: resume qu\'e tanto se mueve la plataforma mientras la base vibra.

\section{Paso 15: C\'omo conectar esto con Simulink y Simscape}

En el enunciado se dice que el modelo realista puede hacerse con bloques est\'andar de Simulink o con componentes de Simscape. En nuestro trabajo:

\begin{itemize}[leftmargin=1.2cm]
\item ya est\'a armado el flujo matem\'atico completo;
\item ya est\'a armado el modelo no lineal y el aproximado;
\item ya est\'a preparada la estructura de comparaci\'on;
\item y se dej\'o una gu\'ia aparte para llevar la planta a una versi\'on f\'isica en Simscape.
\end{itemize}

\subsection{Qu\'e parte ya est\'a lista}

Ya est\'a lista la l\'ogica del proyecto:

\begin{itemize}[leftmargin=1.2cm]
\item modelo realista;
\item modelo aproximado;
\item controlador;
\item referencia;
\item perturbaci\'on;
\item comparaci\'on de resultados.
\end{itemize}

\subsection{Qu\'e parte requiere validaci\'on en tu MATLAB}

La parte que todav\'ia depende de correrlo y revisarlo en MATLAB es la validaci\'on final del diagrama y, si se quiere, el cierre completo del modelo en Simscape.

\section{Paso 16: C\'omo lo debes contar oralmente}

Si te preguntan ``\textit{qu\'e hicieron?}'', una forma muy buena de responder es esta:

\begin{enumerate}[leftmargin=1.2cm]
\item Partimos de un sistema f\'isico realista: una plataforma antivibratoria activa.
\item Obtuvimos un modelo no lineal usando la segunda ley de Newton y un modelo de actuador de primer orden.
\item Lo escribimos en espacio de estado.
\item Linealizamos alrededor del equilibrio de reposo.
\item Obtuvimos una funci\'on de transferencia SISO de tercer orden.
\item Verificamos estabilidad con polos y con Routh-Hurwitz.
\item Dise\~namos un controlador por IMC usando el modelo aproximado.
\item Comparamos el desempe\~no del controlador sobre el modelo aproximado y sobre el modelo realista frente a cambios de referencia y perturbaciones.
\end{enumerate}

\section{Resumen final}

Lo que hicimos, visto con el lenguaje matem\'atico del curso, fue:

\begin{itemize}[leftmargin=1.2cm]
\item comenzar con un sistema no lineal de la forma $\dot{x}=f(x,u)$;
\item escoger variables de estado adecuadas;
\item linealizar alrededor de un punto de equilibrio;
\item obtener la funci\'on de transferencia SISO;
\item analizar estabilidad con polos y Routh-Hurwitz;
\item caracterizar la respuesta temporal;
\item dise\~nar el controlador con IMC;
\item validar el controlador sobre un modelo m\'as realista.
\end{itemize}

Ese es exactamente el hilo l\'ogico que conecta el proyecto con lo que has visto en las diapositivas de estabilidad, respuesta escal\'on, PID e IMC.

\end{document}
